\section{Experiments}
\label{sec:experiments}

\subsection{Dataset}

We evaluate on HotpotQA~\cite{yang2018hotpotqa} in the distractor
configuration, using the 7{,}405-question validation split. Each
example comprises a natural-language question, an answer, an answer
type (``bridge'' or ``comparison''), and ten candidate paragraphs drawn
from Wikipedia. Exactly two of the ten candidates are \emph{gold}
supporting paragraphs; the remaining eight are \emph{distractors}
selected to be topically close to the question. Sentence-level
\texttt{supporting\_facts} annotations pinpoint the specific gold
sentences within the two gold paragraphs. We load the dataset through
the HuggingFace \texttt{hotpotqa/hotpot\_qa} loader and discard the
rare examples whose gold-paragraph count differs from two.

\subsection{Protocol}

We run three seeds, $\text{seed}\in\{0,1,2\}$. For each seed, we
shuffle the validation split with that seed's RNG and take the first
300 questions as the working subset. We then split this subset
100/200 into a training fold (used exclusively to fit
TrustScore-Learned) and a held-out test fold (used to report every
number in the Results section). No trust weight has ever seen a
question on which its Precision@2 is reported; BM25, TF-IDF, SBERT,
Entity-overlap, and TrustScore-Uniform are pointwise deterministic and
involve no training step at all.

The supervised variant fits a scikit-learn logistic regression~\cite{pedregosa2011sklearn}
with $L_2$ regularization over a six-dimensional min-max-normalized
feature vector. Normalization is performed within the candidate pool
of each question, so that per-question score distributions are
comparable.

\subsection{Models and libraries}

SBERT embeddings are produced by
\texttt{sentence-transformers/all-MiniLM-L6-v2}, a compact
sentence-transformer model~\cite{reimers2019sbert}; the model is run
on CPU with normalized output embeddings. BM25 uses the Okapi
variant~\cite{robertson2009bm25} implemented in \texttt{rank\_bm25}.
TF-IDF uses scikit-learn's \texttt{TfidfVectorizer} with default
settings, fit per question on $\{q\}\cup\mathcal{P}$. In the body text
we refer to the sentence-transformer as \emph{all-MiniLM-L6-v2}; the
full org-prefixed identifier above is the canonical HuggingFace hub
path.

End-to-end wall-clock for three seeds is under fifteen minutes on a
single CPU core, comfortably within a laptop budget.

\subsection{Baselines and conditions}

We compare seven conditions, all scored over the same 200-question
held-out fold per seed:

\begin{itemize}
  \item \textbf{Random}: permute the ten candidates with the seed's
  RNG, take the first two. Establishes the floor.
  \item \textbf{BM25}~\cite{robertson2009bm25}: rank by
  $\operatorname{BM25}(q, p_i)$ on a lowercased bag of tokens.
  \item \textbf{TF-IDF cosine}: cosine between TF-IDF vectors of $q$
  and $p_i$, with the vocabulary fit jointly over
  $\{q\}\cup\mathcal{P}$.
  \item \textbf{SBERT}~\cite{reimers2019sbert}: cosine between
  normalized all-MiniLM-L6-v2 embeddings. This is the strongest
  single-signal baseline in our setting.
  \item \textbf{Entity-overlap (ours)}: a standalone diagnostic that
  scores $p_i$ purely by the fraction of question-entity tokens
  present, providing a single-signal reference for the authority
  dimension.
  \item \textbf{TrustScore-Uniform (ablation)}: the uncalibrated
  combination of relevance, length, and consistency defined in the
  Method section.
  \item \textbf{TrustScore-Learned (ours)}: the supervised six-feature
  logistic regression. The proposed method.
\end{itemize}

\subsection{Evaluation metrics}

Given a top-two selection $\hat{\mathcal{S}}$ and the gold set
$\mathcal{S}^{*}=\{i : y_i=1\}$, we report four metrics.
\textbf{Precision@2} is the indicator $\mathbf{1}[\hat{\mathcal{S}}=\mathcal{S}^{*}]$,
averaged over the test fold; it is the strict top-pair metric and our
primary reporting quantity. \textbf{Paragraph F1} is the F1 of
$\hat{\mathcal{S}}$ against $\mathcal{S}^{*}$, more forgiving of
one-correct-one-wrong selections. \textbf{Supporting-fact F1} is the
sentence-level F1 in which the predicted sentence set is the union of
all sentences of the two selected paragraphs; this upper-bounds a
natural downstream pipeline that forwards whole selected paragraphs to
a reader. \textbf{Score gap} is the difference between the mean score
assigned to gold paragraphs and the mean score assigned to
distractors, averaged over the test fold; unlike the ranking metrics,
it is a calibration diagnostic, and a method can rank well (high P@2)
yet remain poorly separated (small gap), which matters when a
downstream module must threshold on the confidence.
